% Chapter 3. Every number is a macro (generated/numbers.tex) or a \needsclaim box.
\chapter{Testbed and Methodology}
\label{ch:methods}

This chapter describes the device, the pipeline, the audio and the protocol
under which every number in Chapters~\ref{ch:cost}--\ref{ch:allocator} was
measured. It closes with the methods results: rules this project adopted
because a measurement went wrong without them. They are reported as results
because each one changed a number.

\section{Device}
\label{sec:device}

All measurements were made on one NVIDIA Jetson Orin Nano Super developer kit
with 8\,GB of memory shared by the CPU and the GPU. The board ran in power mode
MAXN\_SUPER with clocks pinned and was operated headless. Clock state, power
mode and software versions are recorded in every run's header.

The board's six CPU cores are split in two. The language-model server runs on
three cores with every layer offloaded to the GPU. The agent, the voice-activity
detector and both speech recognizers share the other three. Offloading was
audited over the runs made after the server's offload guard was added: none
ran in a server segment that warned of a missing GPU, and the segments' decode
rates had a median of \LlmRate\,ms per token.

\section{Pipeline}
\label{sec:pipeline}

The agent is a cascade built on Pipecat:
\begin{itemize}
\item Silero voice-activity detection, which declares an endpoint after a
  fixed silence;
\item two Whisper engines through faster-whisper with int8 weights:
  \texttt{tiny} decodes partial hypotheses while the user speaks, and
  \texttt{base} decodes the final transcript at the endpoint;
\item a 4-bit Gemma~4~E2B language model served by llama.cpp, with one slot;
\item Piper for synthesis.
\end{itemize}
Each turn writes one record. It holds a timestamp for every stage boundary on a
single clock, every partial and hypothesis decode with its start and end, the
board's power rails, and process memory. Speculation arms add their own
records: drafts, commits, freezes and allocation decisions, each with the
inputs it was taken on.

TTFA is the time from the true end of speech to the first reply audio. The
true end of speech is known to the audio source, not inferred by the pipeline,
so a late endpoint shows up as latency rather than being absorbed into the
definition.

\section{Audio and sets}
\label{sec:sets}

Audio enters through the microphone's delivery path, one file per turn, at
real time. The next file starts only after the reply has finished playing and
a quiet period has passed. Earlier designs used a fixed gap between files and
let a reply overlap the next utterance; the turn gate replaced the fixed gap
(Section~\ref{sec:methods-results}).

The sets:
\begin{itemize}
\item a development set of 16 short spoken assistant queries;
\item from Spoken-MQA \cite{wei2025towards}: two digit splits, a single-step
  split used to select configurations and to test preemption, and the
  multi-step split, derived from GSM8K \cite{cobbe2021gsm8k}. Of the
  multi-step split we use 80 items drawn with a fixed seed and, as a fifth
  length, the 20 longest of them.
\end{itemize}

\section{Protocol}
\label{sec:protocol}

\textbf{Arms and blocking.} Arms are run in blocks with a same-arm washout
turn, or interleaved within one run where the comparison needs it. Every
comparison that carries a claim has a same-day baseline.

\textbf{Pairing and intervals.} Arms are compared on the same items. In the
main comparisons each item's value is its median over three runs, the
difference is taken per item, and the 95\% confidence interval is a
percentile bootstrap over items. Analyses that depart from this state their
design with their result.

\textbf{Validity.} Turns are checked against the file's ground truth. A turn
split by the voice-activity detector, a turn whose endpoint fired after its
final transcript, and a run that fails a registered validity check are
excluded by rules fixed before scoring.

\textbf{The fence.} Between a measured run's start and its teardown, nothing
else runs on the board: no edits, no builds, no tests, no analysis. One breach
is on record. Its run is reported both with and without the affected
repetition, and the headline uses the run without it
(Chapter~\ref{ch:cost}).

\textbf{Pre-registration.} Every prediction, its bar and its scoring rule were
committed to the lab notebook before the run that tests it. A correction is
appended, never edited in. Each pass is scored once, by the rule registered
for it; a prediction tested in more than one pass has one ledger entry per
pass, and a re-scoring or a reconciliation is recorded as such. The full ledger
is in Appendix~\ref{app:prereg}. It records where a prediction was stated
informally, where a bar moved before its data, and where a scorer was written
after the data.

\textbf{Provenance.} Every number in this thesis is produced by a script in
the repository from a log of a measured run. A claim registry pins each
scorer's invocation to the run identifiers the notebook records. Each claim is
one line of that scorer's output, and the thesis cites it through a generated
macro, so no measured value is typed by hand.

\section{Methods results}
\label{sec:methods-results}

Each rule below was adopted after a specific number turned out wrong without
it.

\textbf{Pair within run, on the item.} A difference of medians across
differently sized audio once suggested a large pipeline overhead in the
recognizer. Paired per segment on fixed audio, the overhead disappeared
(\needsclaim{pipeline minus isolation, paired per segment}{2026-09-16 --- STT
inflation attributed}). Every claim-bearing comparison since pairs on the
item; where a design pairs across runs or resamples runs, the result says so.

\textbf{No underpowered null is reported as a null.} The triggers evaluated
in Chapter~\ref{ch:value} have AUCs of \TsemAUC{} (semantic) and \EpaAUC{}
(acoustic) over \TriggerPoints{} decision points from \TriggerUtts{} utterances
of the development set, and \TriggerUnindexed{} points with no utterance index.
Intervals that wide support ``not detectable on this set'', not
``does not predict'', and that is how they are reported.

\textbf{Report marginal rates.} An agreement statistic between two
recognizers once read as perfect because neither had fired; the trigger was
dead on punctuated partials. Every agreement, accuracy or firing statistic is
reported with the base rate it rests on, and a statistic over a constant is
flagged as degenerate.

\textbf{Compare on common support.} When one configuration covers items that
another does not, paired deltas are reported on the items both cover and on
the newly covered items separately. Pooled, one comparison changed sign
(\needsclaim{anticipation lead, pooled vs common support}{Two-engine STT,
live (B1)}).

\textbf{A controller has to vary.} A controller whose output is constant is
reported with its inputs and its decision log, not as a policy. The
allocator's \AllocSpends{} spends in \AllocDecisions{} decisions are reported
that way, with the reason for each decision (Chapter~\ref{ch:allocator}).

\textbf{Derive corpus-dependent constants.} A constant that depends on the
corpus is derived at run time or asserted at start. The feasibility gate's
duration prior is derived at run start as the median duration of the set's
files and recorded in the run header. Constants fixed at registration, such as
the break-even usability in Chapter~\ref{ch:value}, are stated as registered
and are not re-derived per corpus.

\textbf{Wait for the event, not a proxy.} Three runs on one day were voided by
a harness that waited for a count, a silence or a clock instead of the event
it depended on. A reply still playing when the next file began made every
later turn one turn late. The turn gate now waits for the reply to finish.
The abort record is written before the abort.

\textbf{Keep retractions in place.} The notebook holds every withdrawn result
next to its correction:
\begin{itemize}
\item a first scoring of the listener's overlap cost that broke its registered
  population;
\item an estimate of the draft window on mixed clocks;
\item a preemption result voided by its own validity check;
\item an interleaved run whose plan was indexed by turn number, not by file
  position.
\end{itemize}
Chapters~\ref{ch:cost}--\ref{ch:allocator} cite the corrected figures and
name the corrections.
